\documentclass[11pt, a4paper]{scrartcl}
% Packages
\usepackage[margin=1.25in]{geometry}
\usepackage{index}
\usepackage{amsbsy} % Bold math symbols
\makeindex
\usepackage{tcolorbox}
\usepackage{varwidth}
\usepackage{amsmath, amssymb}
\usepackage{esint}
\usepackage{titlesec}
\usepackage{xcolor}
\usepackage{titling}
\usepackage[linktocpage]{hyperref}
\usepackage{pgfplots}
\pgfplotsset{compat=1.18}
\usepackage{caption}
\usepackage{amsthm}
\usepackage{import}
\usepackage{cancel}
\usepackage{caption}
\usepackage{nicematrix}
\usepackage{mathtools}
\usepackage{enumerate}
\usepackage{braket}
\usepackage{graphicx}
\usepackage{booktabs}
\usepackage{makecell}
\usepackage{tikz}
\usepackage[version=4]{mhchem}
\usetikzlibrary{arrows.meta, calc ,positioning, angles,quotes,patterns,patterns.meta,intersections}
\usepackage[italian]{babel}


% Font
%\usepackage{fontspec}
%\setmainfont{Libertinus Serif}

%\usepackage{unicode-math}
%\setmathfont{Libertinus Math}


% To reset footnote numbering each page
\usepackage[perpage]{footmisc}
\usepackage{setspace}
\setstretch{1.2}

% Titles 
\title{Note di\\ \vspace{.15cm} Laboratorio di fisica medica}
\author{}
\date{}


% svolgimento
\newenvironment{svolgimento}{\renewcommand\qedsymbol{$\blacksquare$}\begin{proof}[Svolgimento]}{\end{proof}}

\newtheoremstyle{style}% name of the style to be used
{5pt}% measure of space to leave above the theorem. E.g.: 3pt
{5pt}% measure of space to leave below the theorem. E.g.: 3pt
{\normalfont}% name of font to use in the body of the theorem
%{15pt}% measure of space to indent
{0pt}% measure of space to indent
{\noindent\bfseries}% name of head font
{}% punctuation between head and body
{ }% space after theorem head; " " = normal interword space
{\thmname{#1}\thmnumber{ #2}{\thmnote{ (#3)}.\ }}


\theoremstyle{style}
\newtheorem{esempio}{Esempio}[section]
\newtheorem{definizione}{Definizione}[section]
\newtheorem{prop}{Proposizione}[section]
\newtheorem{teorema}{Teorema}[section]
\newtheorem{lemma}{Lemma}[teorema]
\newtheorem{corollario}{Corollario}[teorema]
\newtheorem{osservazione}{Osservazione}[section]
\newtheorem{assunzione}{Assunzione}[section]
\newtheorem{notazione}{Notazione}[section]
\newtheorem{esercizio}{Esercizio}[section]

%% Generic box
\newtcolorbox{eqbox}[1][]
{
colback=gray!10,
arc=0pt,
boxrule=0pt,
title=#1
}

 \newenvironment{boxenv}[1][]{
    \begin{eqbox}[#1]
    }{
   \end{eqbox}
}

%Captions
\captionsetup[figure]{font=footnotesize,labelfont=footnotesize}
\captionsetup[table]{font=footnotesize,labelfont=footnotesize}
%Titlesec
\titleformat{\section}
{\fontsize{15}{20}\sffamily\scshape}
{\normalfont\color{gray}{\fontsize{15}{20}\selectfont\thesection}}
{0.7em}
{}
\hypersetup{colorlinks,breaklinks, linkcolor=[RGB]{74, 122, 164}}
\definecolor{asdf}{HTML}{4a7aa4}
% Personalizza la formattazione della subsection
\titleformat{\subsection}[block]{\fontsize{13}{20}\bfseries}{\normalfont\thesubsection}{.5em}{}


% Personalizza la formattazione della subsubsection
\titleformat{\subsubsection}[block]{\fontsize{12}{20}\bfseries}{\normalfont\thesubsubsection}{.5em}{}

% Maketitle customization
\renewcommand{\maketitle}{
\begin{center}
{\sffamily
{\fontsize{20}{20}\selectfont\MakeUppercase\thetitle}}

\vspace{0.2in}

{\large\scshape\sffamily\theauthor}
\end{center}
}


% Differenziale 
\newcommand{\dd}{\mathop{}\!\mathrm{d}}




%Evaluate symbol
\DeclareMathOperator{\di}{d\!}
\newcommand*\Eval[3]{\left.#1\right\rvert_{#2}^{#3}}

%%%%%%% Numero delle equazioni in formato a.b
\numberwithin{equation}{subsection}
%%%%%

%%%%%%%%%% Personalizzazione numeri lista
\renewcommand{\theenumi}{(\arabic{enumi})}

%%%% Table of contents

\usepackage[titles]{tocloft}

\renewcommand{\cftdot}{}
\usepackage{titletoc}
%\setcounter{tocdepth}{2}

%%%%%%%%%%%%%%%% Toc style

% Personalizzazione scritta indice


%%%% Lambda bar
\newcommand{\lbar}{{\mkern0.75mu\mathchar '26\mkern -9.75mu\lambda}}

\renewcommand{\arraystretch}{1.3}
\begin{document}
\maketitle
\newpage
\tableofcontents 
\newpage
\section{Rivelatori}

Un rivelatore trasforma l'energia depositata dalla radiazione che lo attraversa in un segnale analogico, elettrico o luminoso, che viene, poi, raccolto ed elaborato da dell'elettronica.
Il tempo dell'interazione $t_s$ \`e dell'ordine di $1$ ns nei rivelatori a gas, mentre \`e dell'ordine di $1$ ps nei rivelatori a stato solido; il tempo di raccolta $t_c$ dell'energia depositata nel rivelatore a seguito dell'interazione \`e $\sim 1$ ms nei rivelatori a gas, mentre \`e $\sim 1$ ns in quelli a stato solido.

Assumendo di avere un rivelatore nel quale viene prodotto un impulso di corrente, che fluisce nel rivelatore per un tempo pari al tempo di raccolta $t_c$, allora la carica totale depositata nel rivelatore \`e
\[
Q = \int _0^{t_c} i(t) \dd t
\] 
\subsection{Modo corrente e modo impulsato}


Se il rate di interazione (numero di eventi rivelati per unit\`a di tempo) \`e molto pi\`u basso rispetto al tempo di raccolta, il rivelatore \`e in grado di rivelare la singola interazione, caratteristica tipica dei rivelatori che operano in \textbf{regime impulsato}.

Quando non si possono distinguere le singole interazioni, il segnale viene raccolto in \textbf{modo corrente}: un picoamperometro campiona il valore di corrente per un tempo $T$, registrando
\[
I(t) = \frac{1}{T} \int_{t-T}^{t} i(t') \dd t'
\] 
Nonostante si perda l'informazione sul singolo evento, il segnale raccolto \`e proporzionale al rate delle interazioni e all'energia media per evento:
\[
I = rQ = r \frac{E}{w} q
\] 
con $r$ rate di eventi, $Q = q E / w$ carica media prodotta per evento, $E$ energia media per evento, $q$ carica elementare e $w$ energia media necessaria per la produzione dell'informazione\footnote{In rivelatori a stato solido, l'informazione \`e data da una coppia elettrone-lacuna; in scintillatori, \`e un fotone luminoso, per esempio.}.

La lettura del segnale in modo impulsato, invece, avviene misurando la differenza di potenziale ai capi della resistenza di un filtro passa alto, dove la capacit\`a $C$ e la resistenza $R$ sono l'equivalente di quelle del rivelatore, dei cavi di connessione e del circuito di misurazione complessivamente.
Si distinguono diversi casi a seconda del tempo caratteristico $\tau  = RC$.
\begin{boxenv}[]
Riprendere da 30:00
\end{boxenv}

















\end{document}
